% BHU PYQ -- Transcribed & Corrected
% Paper: ESSE-21: Geology (Gemology & Earth System Science)
% Exam: B.Sc. (Hons.) Semester II Examination, 2025-26 (NEP)
% Department: Geology

\documentclass[11pt,a4paper]{article}
\usepackage[a4paper,top=2.5cm,bottom=2.5cm,left=2.8cm,right=2.8cm,headheight=14pt]{geometry}
\usepackage{amsmath,amssymb}
\usepackage[T1]{fontenc}
\usepackage[utf8]{inputenc}
\usepackage{lmodern,microtype,fancyhdr,enumitem,hyperref,lastpage,parskip}

\hypersetup{
    colorlinks=true,
    urlcolor=black,
    linkcolor=black,
    pdftitle={ESSE21: Geology -- BHU B.Sc. Sem II 2025-26 (NEP)}
}

\pagestyle{fancy}
\fancyhf{}
\renewcommand{\headrulewidth}{0.4pt}
\renewcommand{\footrulewidth}{0.4pt}
\fancyhead[L]{\small   \textit{Banaras Hindu University}}
\fancyhead[R]{\small   \textit{ESSE-21: Geology (Gemology)}}
\fancyfoot[C]{\small Page  \thepage\ of \pageref{LastPage}}


\newlist{parts}{enumerate}{2}
\setlist[parts,1]{label=(\alph*),leftmargin=*,itemsep=4pt,topsep=2pt}
\setlist[parts,2]{label=(\roman*),leftmargin=*,itemsep=2pt,topsep=2pt}

\newcommand{\pts}[1]{\hfill   \textit{[#1]}}


\begin{document}


\begin{center}
  {\large\bfseries BANARAS HINDU UNIVERSITY}\\[2pt]
  {\normalsize B.Sc. (Hons.) --- Semester II Examination, 2025-26 (NEP)}\\[6pt]
  \rule{0.8\linewidth}{0.5pt}\\[6pt]
  {\large\bfseries GEOLOGY}\\[4pt]
  {\large\bfseries Paper Code: ESSE-21 : Gemology \& Earth System Science}\\[4pt]
  {\normalsize Time: 2:30 Hours\hfill Maximum Marks: 70}\\[2pt]
  {\small REF NO. Conf/Even/N/2025-26/53}
\end{center}

\rule{\linewidth}{0.4pt}

\smallskip



\noindent   \textit{Instructions: Answer   \textbf{five} questions in all including Question No.~1 which is compulsory. All questions carry equal marks (14 marks each).

\bigskip

\centerline{   \textbf{SECTION A}}

\medskip




\noindent   \textbf{Question 1.} What are Precious Gemstones or Maharatnas? Write in detail about Navaratnas or nine precious gemstones with respect to chemical composition, diagnostic physical properties, origin, economic uses, and distribution in India.\pts{14}

\medskip




\noindent   \textbf{Question 2.} Describe the causes of colouration in gemstones in detail.\pts{14}

\medskip




\noindent   \textbf{Question 3.} Briefly describe the following physical properties of minerals with suitable examples: diaphaneity, chatoyancy, asterism, cleavage, parting, tenacity, and specific gravity.\pts{14}

\bigskip
\rule{\linewidth}{0.2pt}
\medskip

\centerline{   \textbf{SECTION B}}

\medskip




\noindent   \textbf{Question 4.} Write chemical composition, diagnostic physical properties, genesis, and economic uses of   \textbf{any two} of the following:\pts{$2  \times 7 = 14$}

\begin{parts}
  \item Igneous rocks
  \item Sedimentary rocks
  \item Metamorphic rocks
  \item Quartz as a semi-precious gemstone
\end{parts}

\medskip




\noindent   \textbf{Question 5.} Write short notes on the following:\pts{$4  \times 3.5 = 14$}

\begin{parts}
  \item Essential parts of a faceted stone
  \item Silver alloys
  \item Isotropic and Anisotropic substances
  \item Classification of diamonds with their specific properties
\end{parts}

\medskip




\noindent   \textbf{Question 6.} Write short notes on   \textbf{any two} methods of gemstone quality enhancement:\pts{$2  \times 7 = 14$}

\begin{parts}
  \item Heat treatment methods
  \item Coating treatment methods
  \item Irradiation treatment methods
  \item Diffusion treatment methods
\end{parts}

\bigskip
\rule{\linewidth}{0.2pt}
\medskip

\centerline{   \textbf{SECTION C (Compulsory)}}

\medskip




\noindent   \textbf{Question 7.} Answer in one/two words or sentences:\pts{14}

\begin{parts}
  \item What are the diagnostic physical characteristics of Corundum?\pts{2}
  \item What are the economic uses of Feldspar?\pts{2}
  \item Write the names of two industrial gemstone synthesis methods.\pts{2}
  \item What is an imitation gemstone?\pts{2}
  \item Name the freshwater and saltwater pearl-bearing bivalve oysters.\pts{2}
  \item Moonstone and sunstone belong to which group of minerals?\pts{1}
  \item In India, the mineral cymophane is named as \underline{\hspace{3cm}}.\pts{1}
  \item Name four crystalline quartzes used as gemstones.\pts{1}
  \item Write the chemical composition of Beryl.\pts{1}
  \item Name the synthetic garnet.\pts{1}
\end{parts}

\vfill
\rule{\linewidth}{0.4pt}

\begin{center}\small
\href{https://bhu-exam.vercel.app/}{Click here to visit our website}\\[4pt]
\href{https://me-aryan.vercel.app/}{Click here to visit our contributor}
\end{center}

\end{document}
